\begin{abstract}
Generating simplified explanations (ELI5-style) requires balancing factual accuracy with accessibility---a challenge that single-pass prompting fails to address adequately, as models must simultaneously handle fact retrieval, logical reasoning, and linguistic simplification in one generation step. We present a multi-agent architecture that decomposes explanation generation through four specialized stages: question breakdown, parallel analysis with retrieval-augmented generation, adaptive synthesis, and creative simplification. This staged approach enables explicit reasoning, external grounding through Wikipedia and curated knowledge bases, and quality-aware content mixing while reducing inference calls by 100$\times$ through batched processing. We evaluate this architecture against baseline single-pass prompting across seven language models (1B-7B parameters) using 30,000 samples and eleven evaluation metrics. Our results reveal a consistent accuracy-quality trade-off: the multi-agent architecture shows a 38.2\% average decline in LLM-judged accuracy but achieves 34.2\% improvement in ROUGE1 and 10.0\% improvement in BERT-F1 scores. Critically, these ROUGE gains are not artifacts of output length: multi-agent explanations average 108 words versus 334 words for baseline (67.6\% shorter) while achieving higher lexical overlap, with a strong negative correlation (r=-0.99) between word count and ROUGE-1 scores. We contribute (1) a novel multi-agent architecture with staged batching that reduces inference calls by 100$\times$, (2) comprehensive empirical evaluation across diverse models and metrics, and (3) the first systematic documentation of the accuracy-quality paradox in multi-agent explanation systems.
\end{abstract}
